\documentclass[conference]{IEEEtran}
\IEEEoverridecommandlockouts
\usepackage{cite}
\usepackage{amsmath,amssymb,amsfonts}
\usepackage{algorithmic}
\usepackage{graphicx}
\usepackage{textcomp}
\usepackage{xcolor}
\usepackage{booktabs}
\usepackage{hyperref}
\usepackage{xurl}

\def\BibTeX{{\rm B\kern-.05em{\sc i\kern-.025em b}\kern-.08em
    T\kern-.1667em\lower.7ex\hbox{E}\kern-.125emX}}

\begin{document}

\title{Proxy Fidelity Under Data Scarcity: Evaluating Nifty TRI as a Stand-In for Unavailable BSE TRI in Indian Open-Source Fintech}

\author{\IEEEauthorblockN{Kushagra Jaiswal}
\IEEEauthorblockA{\textit{NAT\_FUNDS Open Source Project}\\
India\\
}
\and
\IEEEauthorblockN{Harleen Khanuja}
\IEEEauthorblockA{\textit{NAT\_FUNDS Open Source Project}\\
India\\
}
}

\maketitle

\begin{abstract}
Open-source mutual fund analytics platforms in India face a structural data-access barrier: BSE Total Return Index (TRI) time-series are unavailable via any public, programmatically accessible API. Within the NAT\_FUNDS platform, we resolved this by implementing a \emph{data-layer alias}, which substitutes the corresponding NSE Nifty TRI series under BSE benchmark keys in the caching layer. This paper evaluates that architectural decision.

\textbf{Critical scope limitation:} Because BSE TRI itself is unavailable, we cannot measure proxy fidelity against the target series directly. All comparisons use publicly accessible BSE Price Index data (Yahoo Finance) as a surrogate baseline, rather than BSE TRI. Our reported correlations are therefore a \emph{lower bound} on true Nifty-to-BSE-TRI fidelity, since Price indices exclude dividend reinvestment.

Over 120 months (July 2016--June 2026, $n=120$), Pearson correlation between BSE Price indices and Nifty TRI proxies ranged from $r=0.965$ to $r=0.975$ for broad-market pairs (Spearman $\rho=0.951$--$0.964$). A convenience sample of 13 named large-cap ETFs showed mean $\Delta\beta=+0.025$ ($SD=0.014$), which is economically negligible. The observed mean $\Delta\alpha=-1.77\%$ ($SD=0.53\%$) is consistent with the known dividend yield accounting identity (Indian equity market dividend yield $\approx1.5$--$2.0\%$ p.a.) and does not constitute independent validation of the proxy against BSE TRI. Rolling-window robustness analysis shows that $|\Delta\beta|$ remains bounded in magnitude across time ($\approx0.025$--$0.090$), though its sign reverses between earlier and recent periods. Structurally non-equivalent index pairs (e.g., Nifty CPSE vs.\ BSE Select Business Groups, $r=0.656$) are identified as a key failure mode requiring explicit engineering guardrails.
\end{abstract}

\begin{IEEEkeywords}
fintech, open-source, data engineering, proxy substitution, mutual funds, benchmark data, India, TRI
\end{IEEEkeywords}

%% ─── Conflict of Interest ────────────────────────────────────────────────────
\section*{Declaration of Interests}
The authors are the primary developers and maintainers of the NAT\_FUNDS open-source platform evaluated in this paper. This study is a self-evaluation: we designed the architectural decision being assessed and selected the fund sample from our live system cache. We are disclosing this relationship to acknowledge the risk of confirmation bias. A fully independent evaluation would require a separate research team, commercial BSE TRI data (e.g., Bloomberg or CMIE Prowess), and a randomly selected fund sample. Readers should interpret all findings with this context in mind.

%% ─── I. Introduction ─────────────────────────────────────────────────────────
\section{Introduction}

The proliferation of open-source retail investment analytics tools has democratized access to financial data, but developers building such systems in emerging markets routinely encounter severe data asymmetry. High-quality index data that institutional participants obtain via commercial terminals (Bloomberg, Refinitiv) is frequently paywalled, session-gated, or programmatically deprecated for retail and open-source consumers \cite{milan2020}.

In India, this asymmetry has a specific and consequential form tied to a 2018 regulatory mandate. The Securities and Exchange Board of India (SEBI) required all mutual fund schemes to benchmark against Total Return Indices (TRI) rather than Price indices starting from February 2018 \cite{sebi2018}. TRI series incorporate dividend reinvestment and therefore reflect the full economic return of the benchmark, which Price indices systematically understate. The mandate created an environment in which TRI-based benchmarking became a regulatory requirement. However, BSE TRI data remained effectively inaccessible to programmatic consumers.

The NSE maintains a functional Nifty Indices API with full historical TRI series available via documented endpoints. The BSE, by contrast, has deprecated its bulk historical download endpoint for TRI data and provides no publicly documented REST or JSON API for TRI history. This asymmetry places open-source developers building analytics for BSE-benchmarked funds in an untenable position: either omit those funds entirely, revert to the SEBI-disapproved Price index, or substitute a structurally similar NSE series as a proxy.

The NAT\_FUNDS platform chose the third path, implementing what we term a \emph{data-layer alias}: Nifty TRI values stored directly under BSE benchmark keys in the JSON caching layer. This paper quantifies the downstream error introduced into risk metrics (Beta, Alpha) by this substitution. Its contributions are: \textbf{(1)} a documented, reproducible methodology for bounding proxy fidelity under data unavailability; \textbf{(2)} empirical metric error estimates for 13 named large-cap ETFs; \textbf{(3)} identification of a structural failure mode for thematically non-equivalent index pairs; and \textbf{(4)} explicit documentation of the fundamental limits of this evaluation.

%% ─── II. Related Work ────────────────────────────────────────────────────────
\section{Related Work}

\subsection{Tracking Error and Benchmark Fidelity}
The relationship between index-tracking instruments and their benchmarks has been studied extensively in developed markets. Rompotis \cite{rompotis2011} documents predictable patterns in ETF tracking error and shows that it is driven by liquidity, expense ratios, and replication methodology. This body of work focuses on fund-to-benchmark deviation; our paper addresses a distinct and less-studied problem: benchmark-to-benchmark deviation caused by data substitution at the infrastructure layer.

\subsection{Data Accessibility and Emerging Market Fintech}
Milan and Treré \cite{milan2020} identify \emph{data universalism} (the implicit assumption that data infrastructure is globally uniform) as a systemic bias in technology research. The BSE API deprecation is a concrete instance of this problem. It creates a structural disadvantage for open-source developers relative to institutional participants, and forces engineering workarounds whose downstream effects are rarely studied or disclosed. To our knowledge, no prior work has empirically evaluated the metric-level consequences of such workarounds in the Indian mutual fund context.

\subsection{Regulatory Context: SEBI's TRI Mandate}
SEBI's 2018 circular mandating TRI-based benchmarking \cite{sebi2018} established the regulatory reason this problem matters: Alpha and return comparisons computed against Price indices systematically overstate fund performance by the dividend yield ($\approx1.5$--$2.0\%$ p.a.\ for Indian equities). The mandate thus made access to TRI data a compliance-adjacent engineering requirement for any platform reporting fund analytics to retail investors.

%% ─── III. System Architecture ────────────────────────────────────────────────
\section{System Architecture: The Data-Layer Alias}

When developing NAT\_FUNDS, we encountered the BSE API deprecation described above. Rather than implementing proxy logic at the metric computation stage---which would require per-fund conditional branching throughout the codebase---we implemented the substitution at the data-persistence layer.

Nifty TRI values are stored directly under BSE benchmark keys in the platform's JSON cache. A system query for \texttt{BSE SENSEX TRI} returns an identical daily NAV array to \texttt{Nifty 50 TRI}; similarly, \texttt{BSE 100 TRI} returns Nifty 100 TRI data, and so on. This design achieves architectural uniformity: all downstream metric calculations (Alpha, Beta, Sharpe, Sortino) operate on the aliased data via identical code paths, without knowledge of the substitution.

This design has one important consequence: because the substitution occurs silently at the data layer, end users of the platform require explicit disclosure that metrics labeled against a BSE index are computed from a Nifty proxy. Based on the findings of this evaluation, the NAT\_FUNDS platform deployed a transparent user-interface disclosure layer, specifically an inline tooltip on all affected fund detail pages, to accurately reflect the proxy substitution to retail investors. The engineering recommendation of the present study is that any deployment of this data-layer alias pattern must include such disclosure to prevent retail harm.

%% ─── IV. Empirical Methodology ───────────────────────────────────────────────
\section{Empirical Methodology}

The empirical evaluation was conducted in four distinct stages to progress from architectural discovery to statistical validation and synthesis.

\subsection{Stage 0: Architectural Discovery and Index Correlation}
\label{sec:stage0}
The target of this evaluation (the error between Nifty TRI and BSE TRI) cannot be measured directly, because BSE TRI is the series whose inaccessibility motivated the substitution. We established an independent baseline using BSE Price Index data from Yahoo Finance (\texttt{\^{}BSESN}, \texttt{BSE-100.BO}, \texttt{BSE-200.BO}, \texttt{BSE-500.BO}). These are the only publicly accessible, long-run BSE series. All reported correlations and metric errors are therefore relative to BSE Price, not BSE TRI. Since Price indices omit dividend reinvestment, the true Nifty-to-BSE-TRI correlation is expected to be \emph{higher} than the values measured here.

Our dataset spans 120 months of monthly returns (July 2016--June 2026, $n=120$), determined by the earliest date of continuous BSE index availability in Yahoo Finance. Pearson correlation coefficients and 95\% confidence intervals were computed using 10{,}000 bootstrap resamples with the percentile method (NumPy \texttt{default\_rng}, seed=42). Spearman rank correlation was computed as a non-parametric robustness check.

\subsection{Stage 1: Metric-Level Error Quantification}
\label{sec:stage1}
To quantify downstream impact, we extracted a convenience sample of 13 BSE-benchmarked mutual fund ETFs from the NAT\_FUNDS live application cache (all fund identifiers are listed in Appendix~\ref{app:funds}). For each fund, we aligned 36 months of monthly returns across three series: Fund NAV, Proxy Benchmark (Nifty TRI from cache), and True Baseline (BSE Price from Yahoo Finance). Three-year Beta was computed as $\hat{\beta} = \text{Cov}(r_f, r_b) / \text{Var}(r_b)$. Annualized Alpha was computed as $(1 + \alpha_\text{monthly})^{12} - 1$, where $\alpha_\text{monthly} = \bar{r}_f - r_f^0 - \hat{\beta}(\bar{r}_b - r_f^0)$ and risk-free rate $r_f^0 = 6.0\%$ p.a.\ (approximate Indian 91-day T-bill yield over the period, compounded monthly). $\Delta\beta = \beta_\text{proxy} - \beta_\text{true}$ and $\Delta\alpha = \alpha_\text{proxy} - \alpha_\text{true}$. Statistical significance was assessed with a two-tailed paired $t$-test.

\subsection{Stage 2: Statistical Validation (Rolling-Window Robustness)}
\label{sec:stage2}
To assess whether the static 36-month results from Stage 1 are representative, we performed a rolling-window analysis. For each fund, $\Delta\beta$ and $\Delta\alpha$ were computed for every contiguous 36-month window available in the full dataset (July 2016--June 2026), advancing by one month per step. For the 13 funds in our sample this produced 849 total (fund, window) observations. The mean, standard deviation, and inter-quartile range (IQR) of $\Delta\beta$ and $\Delta\alpha$ across all windows measure temporal stability of the proxy effect.

\subsection{Stage 3: Synthesis and Reproducibility}
\label{sec:stage3}
The final stage synthesized findings from Stages 0--2 into a unified evaluation of proxy fidelity. The core platform codebase is available at the project repository root (\url{https://github.com/kush0712/NAT_FUNDS}), while all Python verification scripts, raw data CSVs, and output matrices are isolated in the \texttt{research/} sub-directory. Within this directory, scripts are strictly separated into two folders: \texttt{phase0/} (for index correlation analysis) and \texttt{phase1/} (for metric-level impact and rolling-window robustness). Each contains its own \texttt{data/} input and \texttt{outputs/} results sub-folders. Dependency versions are documented in \texttt{research/requirements.txt}, and a step-by-step reproduction guide is provided in \texttt{research/README.md}.

%% ─── V. Results & Discussion ─────────────────────────────────────────────────
\section{Results \& Discussion}

\subsection{Broad Market Correlation}
Broad-market BSE Price indices exhibit high co-movement with their Nifty TRI proxies (Table~\ref{tab:correlation}). Pearson $r$ ranged from $0.965$ to $0.975$ across the four broad-market pairs; all Spearman $\rho$ values fell within $0.013$ of the corresponding Pearson values, confirming the relationship is not driven by outliers. All 95\% bootstrap confidence intervals are bounded above $0.934$.

\begin{table}[htbp]
\caption{BSE Price vs.\ Nifty TRI: 120-Month Pearson and Spearman Correlation\\(July 2016--June 2026; bootstrap 95\% CI, 10{,}000 resamples, percentile method)}
\begin{center}
\begin{tabular}{lccc}
\toprule
\textbf{Pair (BSE Price $\leftrightarrow$ Nifty TRI)} & \textbf{$r$} & \textbf{95\% CI} & \textbf{$\rho$} \\
\midrule
BSE SENSEX $\leftrightarrow$ Nifty 50       & 0.965 & [0.940, 0.980] & 0.951 \\
BSE 100 $\leftrightarrow$ Nifty 100         & 0.973 & [0.955, 0.985] & 0.964 \\
BSE 200 $\leftrightarrow$ Nifty 200         & 0.973 & [0.956, 0.984] & 0.964 \\
BSE 500 $\leftrightarrow$ Nifty 500         & 0.975 & [0.958, 0.985] & 0.964 \\
\bottomrule
\end{tabular}
\label{tab:correlation}
\end{center}
\end{table}

As noted in Section~\ref{sec:stage0}, these values are lower bounds: BSE TRI series (unavailable) include dividend reinvestment and are expected to co-move more tightly with Nifty TRI than BSE Price does.

\subsection{Thematic Divergence Risk}
Structurally non-equivalent index substitutions produce substantially lower correlations (Table~\ref{tab:thematic}). The Nifty CPSE index (Public Sector Enterprises, government-owned companies) substituted for BSE Select Business Groups (family-controlled conglomerates) yields $r=0.656$ and $\rho=0.674$. This reflects genuine compositional divergence. Nifty 200 Quality 30 (quality-factor weighted) vs.\ BSE Quality TRI yields $r=0.889$. These are not measurement artifacts; they reflect fundamentally different sector exposures and return drivers. Automated wholesale proxy mapping across all BSE indices, without constituent-level equivalence verification, is a significant engineering risk.

\begin{table}[htbp]
\caption{Non-Equivalent Index Pairs: 120-Month Correlation}
\begin{center}
\begin{tabular}{lccc}
\toprule
\textbf{Pair} & \textbf{$r$} & \textbf{95\% CI} & \textbf{$\rho$} \\
\midrule
BSE Sel.\ Bus.\ Groups $\leftrightarrow$ Nifty CPSE & 0.656 & [0.560, 0.735] & 0.674 \\
BSE Quality TRI $\leftrightarrow$ Nifty 200 Qual.\ 30 & 0.889 & [0.833, 0.927] & 0.838 \\
\bottomrule
\end{tabular}
\label{tab:thematic}
\end{center}
\end{table}

\subsection{Beta Stability}
Across the 13-fund convenience sample, the proxy introduces a mean $\Delta\beta=+0.025$ ($SD=0.014$). A two-tailed paired $t$-test yields $t(12)=6.24$, $p<0.001$, confirming that the shift is statistically non-zero. However, an absolute Beta shift of $0.025$ is economically negligible for retail fund evaluation: at a market return of $12\%$ p.a., this corresponds to an additional $0.30\%$ performance-attribution error. From a fitness-for-purpose standpoint, the proxy preserves Beta integrity for broad-market funds.

\subsection{Proxy vs.\ BSE Price Alpha Differential}
Mean $\Delta\alpha=-1.77\%$ ($SD=0.53\%$). This value is most accurately characterized as a \emph{confirmation of a known accounting identity} rather than as independent validation of proxy quality.

BSE Price indices exclude dividend reinvestment. Computing a fund's Alpha against a Price benchmark therefore artificially inflates apparent manager outperformance by approximately the market dividend yield. The Indian equity market has sustained a historical dividend yield of $\approx1.5$--$2.0\%$ p.a., which fully accounts for the observed $\Delta\alpha$. That is: $\Delta\alpha \approx -(\text{dividend yield}) \times 1 \approx -1.77\%$, consistent with basic index-return decomposition.

The finding does establish that the Nifty TRI proxy \emph{avoids} this Price-index bias for the 13 funds in our sample, which is a genuine benefit. Alpha computed against Nifty TRI is a more honest representation of manager skill than Alpha computed against a BSE Price benchmark. What the finding does not establish is whether Nifty TRI is a good proxy for BSE TRI specifically, as that question requires data we do not have.

\subsection{Rolling-Window Robustness}
Rolling-window analysis (849 total (fund, window) observations; Section~\ref{sec:stage2}) reveals that the static 36-month results are \emph{not} representative of the full distribution and should be interpreted with caution.

Across all 36-month windows from July 2016 to June 2026, mean $\Delta\beta = -0.088$ ($SD=0.046$; IQR $[-0.112, -0.064]$). The static window (most recent 36 months) yielded $\Delta\beta = +0.025$, which has a different sign than the cross-window mean. This sign reversal reflects a genuine change in the relative volatility structures of the BSE and NSE indices over time, and is itself an important finding: the proxy's Beta effect is not stable in direction across the full history. It appears positive in recent years and negative in earlier periods, with the magnitude consistent across time ($|\Delta\beta| \approx 0.025$--$0.090$).

Mean $\Delta\alpha = -0.51\%$ ($SD=1.19\%$; IQR $[-1.27\%, +0.12\%]$) across all windows. The wide IQR indicates substantial variation in the Alpha differential across time periods. In some windows the proxy actually increases apparent Alpha, and in others it decreases it. This is consistent with year-to-year variation in the Indian equity dividend yield and changing relative performance between BSE and NSE market segments.

These rolling-window results materially strengthen the case for the threats identified in Section~\ref{sec:threats}. The static 36-month window reported in Sections~V-C and V-D reflects the most recent period, which may not be typical. Researchers and practitioners adopting this proxy should account for the sign variability in $\Delta\beta$ when interpreting historical performance.

%% ─── VI. Threats to Validity ─────────────────────────────────────────────────
\section{Threats to Validity}
\label{sec:threats}
\subsection{Self-Evaluation (Internal Validity)}
The most significant threat is that the authors are evaluating their own architectural decision. We mitigated this by pre-writing analysis scripts before the final results were tabulated and by reporting the full result distribution including failure cases (Section~V-B). However, unconscious confirmation bias in metric selection, baseline choice, and result framing cannot be ruled out. An independent evaluation by a separate team with access to commercial BSE TRI data would provide substantially stronger evidence.

\subsection{Convenience Sample and Selection Bias (External Validity)}
The 13 funds were drawn from the NAT\_FUNDS live cache---a system that already implements the substitution we are evaluating. Funds that the platform handles without error are therefore overrepresented. The sample consists predominantly of large-cap BSE SENSEX ETFs from major AMCs. These funds are expected to have the highest benchmark fidelity. Results for mid-cap, small-cap, or sector-specific BSE-benchmarked funds may differ and were not assessed. No power analysis was performed to justify $n=13$.

\subsection{Proxy-vs-Proxy Comparison (Construct Validity)}
The fundamental epistemological limit of this study is that the target series (BSE TRI) is identical to the series whose inaccessibility motivated the substitution. All reported correlations and metric errors are relative to BSE Price, not BSE TRI. Our results bound proxy fidelity \emph{from below} but provide no upper bound. The true fidelity of Nifty TRI as a stand-in for BSE TRI could be substantially higher or lower than the values reported here; we simply cannot measure it with available public data.

\subsection{Single Time Window (Temporal Validity)}
The analysis covers July 2016--June 2026 only. The rolling-window check (Section~V-E) provides internal evidence of stability across sub-windows, but the full period covers a single sustained bull market with two notable disruptions (2020 COVID crash, 2022 rate-cycle shock). The proxy relationship during structural market regime changes or periods of differential BSE/NSE liquidity has not been evaluated.

%% ─── VII. Conclusion ─────────────────────────────────────────────────────────
\section{Conclusion}

The data-layer alias implemented in NAT\_FUNDS (which substitutes Nifty TRI for unavailable BSE TRI) introduces a statistically significant but economically small Beta shift in the most recent 36-month window ($\Delta\beta\approx+0.025$), and avoids the Price-vs-TRI dividend accounting error that would arise from using the only publicly available BSE alternative. Structurally non-equivalent thematic index pairs (Nifty CPSE for BSE Select Business Groups, $r=0.656$) represent a significant failure mode requiring explicit engineering guardrails, such as a curated equivalence map validated at the constituent level.

However, rolling-window analysis reveals that the direction of $\Delta\beta$ is not stable across the full 2016--2026 history (cross-window mean $-0.088$, IQR $[-0.112, -0.064]$), with the static window showing a different sign than the historical average. This nuance implies that the proxy effect is highly sensitive to the evaluation period, and should not be treated as a static constant.

Four limitations materially bound these conclusions: (1) all comparisons are proxy-vs-price, not proxy-vs-target; (2) the sample is a convenience selection from the authors' own system; (3) this is a self-evaluation; and (4) the static window is not representative of the full cross-window distribution. The findings are best understood as \emph{engineering necessity bounds}: The substitution performs within acceptable limits given what can be measured, but the central question of whether Nifty TRI faithfully tracks BSE TRI remains open pending commercial BSE TRI data.

The primary actionable recommendation for practitioners adopting this pattern is transparent UI disclosure: any metric computed against a substitute benchmark should be annotated as such at the point of display, with the correlation of the substitute documented for the user.

%% ─── References ──────────────────────────────────────────────────────────────
\begin{thebibliography}{00}

\bibitem{milan2020}
S. Milan and E. Treré, ``Big Data from the South(s): Beyond Data Universalism,''
\textit{Television \& New Media}, vol.~20, no.~4, pp.~319--335, 2019.

\bibitem{rompotis2011}
G. Rompotis, ``Predictable patterns in ETFs' return and tracking error,''
\textit{Studies in Economics and Finance}, vol.~28, no.~1, pp.~14--35, 2011.

\bibitem{sebi2018}
Securities and Exchange Board of India, ``Circular: Performance Benchmarking of
Mutual Fund Schemes against Total Return Index (TRI),''
SEBI/HO/IMD/DF3/CIR/P/2018/69, June 2018. [Online]. Available:
\url{https://www.sebi.gov.in/legal/circulars/jun-2018/circular-on-performance-benchmarking-of-mutual-fund-schemes-against-total-return-index-tri-_39078.html}

\end{thebibliography}

%% ─── Appendix ────────────────────────────────────────────────────────────────
\appendix
\section{Fund Sample}
\label{app:funds}

Table~\ref{tab:funds} lists all 13 fund schemes in the convenience sample, with their BSE benchmark assignments as recorded in the NAT\_FUNDS \texttt{benchmark-data.json} file.

\begin{table}[htbp]
\caption{Convenience Sample: Fund Scheme Names and BSE Benchmarks}
\begin{center}
\begin{tabular}{p{5.2cm}c}
\toprule
\textbf{Scheme Name} & \textbf{Benchmark} \\
\midrule
SBI BSE 100 ETF                       & BSE 100 TRI \\
LIC MF BSE Sensex ETF                 & BSE SENSEX TRI \\
Aditya Birla Sun Life BSE Sensex ETF  & BSE SENSEX TRI \\
SBI BSE SENSEX ETF                    & BSE SENSEX TRI \\
UTI BSE Sensex ETF                    & BSE SENSEX TRI \\
BANDHAN BSE Sensex ETF                & BSE SENSEX TRI \\
ICICI Prudential BSE 500 ETF          & BSE 500 TRI \\
DSP BSE Sensex ETF                    & BSE SENSEX TRI \\
Nippon India ETF BSE Sensex           & BSE SENSEX TRI \\
ICICI Prudential BSE Sensex ETF       & BSE SENSEX TRI \\
Mirae Asset BSE Sensex ETF            & BSE SENSEX TRI \\
Kotak BSE Sensex ETF                  & BSE SENSEX TRI \\
Axis BSE Sensex ETF                   & BSE SENSEX TRI \\
\bottomrule
\end{tabular}
\label{tab:funds}
\end{center}
\end{table}

\end{document}
